\documentclass[10pt]{article}

% --------------------------------------------------------
% PACOTES
% --------------------------------------------------------
\usepackage[brazil]{babel}     
%\usepackage[latin1]{inputenc}
%\usepackage{ae}
%\usepackage{epsfig}
%\usepackage{amsmath,amssymb}
%\usepackage{multirow}
%\usepackage{url}
%\usepackage[final]{pdfpages}
\usepackage[utf8]{inputenc}


% Pacotes de citações
% ---
\usepackage[brazilian,hyperpageref]{backref}	 % Paginas com as citações na bibl
%\usepackage[alf]{abntex2cite}	% Citações padrão ABNT

\usepackage{natbib}
\usepackage{graphicx}
\usepackage{indentfirst}		
\usepackage{lscape}
\usepackage{array}
\usepackage{scalefnt}
\usepackage{longtable}
\usepackage{tabularx}
\usepackage{ragged2e}
\usepackage[table]{xcolor}

% --------------------------------------------------------
% INFORMAÇÕES DE DADOS PARA PDF FINAL
% --------------------------------------------------------

\title{Título do Projeto de Pesquisa}
\author{Autor}
\date{Ano}


% --------------------------------------------------------

% --------------------------------------------------------
% ESPAÇAMENTOS ENTRE LINHAS E PARÁGRAFOS
% --------------------------------------------------------

% O tamanho do parágrafo é dado por:
\setlength{\parindent}{1.3cm}

% Controle do espaçamento entre um parágrafo e outro:
\setlength{\parskip}{0.2cm}

% --------------------------------------------------------
% INÍCIO DO DOCUMENTO
% --------------------------------------------------------

\begin{document}

% Seleciona o idioma do documento (conforme pacotes do babel)
\selectlanguage{brazil}

% Retira espaço extra obsoleto entre as frases.
\frenchspacing


% --------------------------------------------------------
% ELEMENTOS PRÉ-TEXTUAIS
% --------------------------------------------------------

\input{capa.tex} % ALTERAR A CAPA NESSE ARQUIVO .TEX

% --------------------------------------------------------
% ELEMENTOS TEXTUAIS
% --------------------------------------------------------

\input{chapters/Introducao}

\input{chapters/Objetivo}

\input{chapters/Metodologia}

\input{chapters/Cronograma}

% --------------------------------------------------------
% REFERÊNCIAS BIBLIOGRÁFICAS
% --------------------------------------------------------

\bibliographystyle{unsrt}
%\bibliographystyle{abnt-alf}
\bibliography{references}

% --------------------------------------------------------
% FINAL DO DOCUMENTO
% --------------------------------------------------------

\end{document}
